\documentclass[12pt,a4paper]{article}

\usepackage[utf8]{inputenc}
\usepackage[T1]{fontenc}
\usepackage[spanish]{babel}
\usepackage{amsmath,amssymb,amsfonts}
\usepackage{graphicx}
\usepackage{booktabs}
\usepackage{hyperref}
\usepackage{natbib}
\usepackage{geometry}
\geometry{margin=2.5cm}

\title{Modelo H\'{i}brido de Presupuesto Energ\'{e}tico Din\'{a}mico y Red Neuronal Artificial para la Estimaci\'{o}n de Biomasa en Larvas de \textit{Hermetia illucens} a partir de la Din\'{a}mica Respiratoria de CO$_2$}
\author{John Jairo Leal G\'{o}mez \\ Universidad Nacional de Colombia, Palmira}
\date{\today}

\begin{document}

\maketitle

\begin{abstract}
\noindent La mosca soldado negra (\textit{Hermetia illucens}, BSF) ha emergido como un agente prometedor para la valorizaci\'{o}n de residuos org\'{a}nicos mediante bioconversi\'{o}n. El monitoreo en tiempo real de la biomasa larvaria durante la crianza es cr\'{i}tico para la optimizaci\'{o}n de procesos, sin embargo, la medici\'{o}n directa sigue siendo destructiva y laboriosa. Este estudio presenta un marco de modelado h\'{i}brido novedoso que combina un modelo mecanicista de Presupuesto Energ\'{e}tico Din\'{a}mico (DEB) con una Red Neuronal Artificial B\'{a}sica (RNA-B) para estimar la biomasa larvaria a partir de mediciones continuas de CO$_2$ respiratorio. El componente DEB proporciona restricciones fisiol\'{o}gicamente informadas sobre la din\'{a}mica de crecimiento, mientras que la RNA-B corrige las desviaciones espec\'{i}ficas del sustrato y la variabilidad ambiental no capturada por la formulaci\'{o}n mecanicista. El modelo fue validado contra experimentos de laboratorio controlados (n=5 r\'{e}plicas, 50 mediciones temporales cada una) bajo diferentes calidades de sustrato. Los resultados demuestran que el enfoque h\'{i}brido supera tanto al DEB independiente (reducci\'{o}n del AME: 34\%) como a la RNA-B pura (reducci\'{o}n del error de generalizaci\'{o}n: 28\%), logrando un error absoluto medio de 2,1 mg en la predicci\'{o}n de biomasa seca. El marco ofrece una soluci\'{o}n escalable y no destructiva para el monitoreo de biomasa en instalaciones industriales de crianza de BSF, con implicaciones para la agricultura de precisi\'{o}n con insectos y la contabilidad de emisiones de gases de efecto invernadero.
\end{abstract}

\noindent\textbf{Palabras clave:} Mosca soldado negra, Presupuesto Energ\'{e}tico Din\'{a}mico, Red Neuronal Artificial, Bioconversi\'{o}n, Monitoreo de CO$_2$, Agricultura de precisi\'{o}n con insectos

\section{Introducci\'{o}n}

El imperativo global de transitar hacia bioeconom\'{i}as circulares ha intensificado la investigaci\'{o}n en sistemas de valorizaci\'{o}n biol\'{o}gica de residuos \citep{goldBiowasteTreatmentBlack2020}. Entre los agentes biol\'{o}gicos, la mosca soldado negra (\textit{Hermetia illucens}, Diptera: Stratiomyidae) ha recibido atenci\'{o}n excepcional debido a su capacidad para convertir diversos sustratos org\'{a}nicos en biomasa rica en prote\'{i}nas y l\'{i}pidos, adecuada para alimentaci\'{o}n animal, biodiesel y fertilizantes org\'{a}nicos \citep{makkarStateoftheartUseInsects2014, liSimultaneousUtilizationGlucose2015}.

A pesar de la r\'{a}pida adopci\'{o}n industrial, la crianza de BSF enfrenta desaf\'{i}os cr\'{i}ticos de control de procesos. La biomasa larvaria -- el principal producto econ\'{o}mico -- se cuantifica tradicionalmente mediante muestreo destructivo, lo que interrumpe los ciclos de crianza, introduce sesgos de muestreo y excluye la toma de decisiones en tiempo real \citep{dienerConversionOrganicMaterial2011}. Las alternativas no destructivas, como la estimaci\'{o}n de biomasa basada en im\'{a}genes, sufren de oclusi\'{o}n por sustrato y requieren calibraci\'{o}n compleja \citep{meneguzEffectRearingSubstrate2018}.

Una alternativa emergente aprovecha el acoplamiento metab\'{o}lico fundamental entre crecimiento y respiraci\'{o}n. El crecimiento de insectos est\'{a} intr\'{i}nsecamente ligado a la producci\'{o}n de CO$_2$ a trav\'{e}s del metabolismo de mantenimiento, los costos de bios\'{i}ntesis y la din\'{a}mica de almacenamiento de l\'{i}pidos \citep{arreseInsectFatBody2010}. \citet{eriksenDynamicModellingFeed2022} demostr\'{o} que un modelo de Presupuesto Energ\'{e}tico Din\'{a}mico (DEB) puede predecir con precisi\'{o}n el crecimiento y las emisiones de CO$_2$ de BSF a trav\'{e}s de diferentes calidades de sustrato. Sin embargo, los modelos DEB requieren par\'{a}metros espec\'{i}ficos de especie que son dif\'{i}ciles de estimar bajo condiciones variables de campo, y su estructura determinista no puede capturar las desviaciones metab\'{o}licas espec\'{i}ficas del sustrato \citep{padmanabhaComprehensiveDynamicGrowth2020}.

Concurrentemente, las Redes Neuronales Artificiales (RNA) han demostrado ser efectivas en tareas de predicci\'{o}n agroindustrial, particularmente cuando las relaciones mecanicistas est\'{a}n parcialmente oscurecidas por la variabilidad ambiental \citep{kamilarisDeepLearningAgriculture2018}. Sin embargo, los enfoques de RNA pura carecen de interpretabilidad fisiol\'{o}gica, requieren datos extensivos de entrenamiento y pueden fallar en generalizar m\'{a}s all\'{a} de las condiciones experimentales observadas \citep{lecunDeepLearning2015}.

Este estudio aborda estas limitaciones a trav\'{e}s de un \textbf{marco de modelado h\'{i}brido} que integra:
\begin{enumerate}
    \item Un \textit{modelo DEB simplificado} que proporciona restricciones de crecimiento informadas fisiol\'{o}gicamente y predicciones de flujos de CO$_2$;
    \item Una \textit{Red Neuronal Artificial B\'{a}sica (RNA-B)} con arquitectura m\'{i}nima, entrenada para corregir los residuos del DEB utilizando trayectorias de CO$_2$ espec\'{i}ficas del sustrato.
\end{enumerate}

El enfoque h\'{i}brido busca combinar la robustez extrapolativa del modelado mecanicista con la flexibilidad adaptativa del aprendizaje autom\'{a}tico, produciendo una herramienta escalable de estimaci\'{o}n de biomasa para instalaciones industriales de BSF.

\section{Materiales y M\'{e}todos}

\subsection{Marco Te\'{o}rico}

\subsubsection{Fundamento del Presupuesto Energ\'{e}tico Din\'{a}mico}

Siguiendo a \citet{eriksenDynamicModellingFeed2022}, la biomasa larvaria se compartimenta en biomasa estructural ($B$, mg) y l\'{i}pidos de almacenamiento ($L$, mg), con un pool de asimilados ($A$) en cuasi-estado estacionario. Las ecuaciones gobernantes son:

\begin{align}
    \frac{dA}{dt} &= 0 = r_A - \left(r_B + r_L + r_{CO_2,m} + r_{CO_2,B} + r_{CO_2,L}\right) \label{eq:assimilate}\\
    \frac{dB}{dt} &= r_B = \mu_B \cdot B \label{eq:structural}\\
    \frac{dL_{tot}}{dt} &= r_L + \delta_{L,B} \, r_B \label{eq:lipid}
\end{align}

donde $r_A$ es la tasa de asimilaci\'{o}n, $r_B$ y $r_L$ son las tasas de s\'{i}ntesis de biomasa estructural y l\'{i}pidos, y $r_{CO_2,\{m,B,L\}}$ representan la producci\'{o}n de CO$_2$ por mantenimiento, s\'{i}ntesis estructural y s\'{i}ntesis de l\'{i}pidos, respectivamente. El par\'{a}metro $\delta_{L,B} = 0,1$ denota la fracci\'{o}n de l\'{i}pidos estructurales \citep{eriksenDynamicModellingFeed2022}.

La tasa espec\'{i}fica de crecimiento $\mu_B$ incorpora una regulaci\'{o}n log\'{i}stica:

\begin{equation}
    \mu_B = \left(\frac{a - m}{1 + Y_B}\right)\left(1 - \left(\frac{B}{B_{max}}\right)^\beta\right) \label{eq:muB}
\end{equation}

donde $a$ es la tasa espec\'{i}fica de asimilaci\'{o}n, $m$ la tasa de mantenimiento, $Y_B$ el costo de bios\'{i}ntesis y $\beta$ el coeficiente de regulaci\'{o}n del crecimiento.

\subsubsection{Producci\'{o}n de CO$_2$ como Observable}

La producci\'{o}n total de CO$_2$ integra tres flujos metab\'{o}licos:

\begin{equation}
    r_{CO_2} = \underbrace{m \cdot B}_{\text{mantenimiento}} + \underbrace{Y_B \cdot r_B}_{\text{s\'{i}ntesis estructural}} + \underbrace{Y_L \cdot r_L}_{\text{s\'{i}ntesis de l\'{i}pidos}} \label{eq:co2_total}
\end{equation}

donde $Y_L = 0,42$ es el costo de s\'{i}ntesis de l\'{i}pidos \citep{kokPreliminaryProjectDesign2021}. En ambientes controlados con temperatura y humedad constantes, $r_{CO_2}$ se convierte en un proxy directo de la actividad metab\'{o}lica y, por extensi\'{o}n, de la acumulaci\'{o}n de biomasa.

\subsubsection{Estructura de Correcci\'{o}n de la Red Neuronal}

La RNA-B se formula como una red feedforward m\'{i}nima con una sola capa oculta, dise\~{n}ada para prevenir el sobreajuste dada la replicaci\'{o}n experimental limitada. La arquitectura es:

\begin{equation}
    \hat{B}_{RNA}(t) = W^{(2)} \cdot \sigma\left(W^{(1)} \cdot \mathbf{x}(t) + \mathbf{b}^{(1)}\right) + b^{(2)} \label{eq:rna}
\end{equation}

donde el vector de entrada $\mathbf{x}(t) = [\Delta C(t), \tau_{50\%}, A_{cum}]$ comprende:
\begin{itemize}
    \item $\Delta C(t) = C(t) - C_0$: exceso acumulado de CO$_2$ (ppm);
    \item $\tau_{50\%}$: tiempo para alcanzar el 50\% de la saturaci\'{o}n de CO$_2$ (d\'{i}a);
    \item $A_{cum}$: \'{a}rea integrada bajo la curva de CO$_2$ (ppm$\cdot$d\'{i}a).
\end{itemize}

La capa oculta emplea activaci\'{o}n ReLU: $\sigma(z) = \max(0, z)$. La capa de salida usa activaci\'{o}n lineal para predicci\'{o}n de biomasa sin cota superior.

\subsection{Integraci\'{o}n H\'{i}brida}

La predicci\'{o}n h\'{i}brida combina la extrapolaci\'{o}n DEB con la correcci\'{o}n residual de la RNA-B:

\begin{equation}
    \hat{B}_{hibrido}(t) = \underbrace{B_{DEB}(t)}_{\text{base mecanicista}} + \underbrace{\gamma(t) \cdot \hat{B}_{RNA}(t)}_{\text{correcci\'{o}n adaptativa}} \label{eq:hybrid}
\end{equation}

donde $\gamma(t) \in [0,1]$ es un peso de confianza dependiente del tiempo, determinado por la discrepancia entre las trayectorias de CO$_2$ observadas y predichas por DEB:

\begin{equation}
    \gamma(t) = \text{sigmoid}\left(\kappa \cdot |r_{CO_2,obs}(t) - r_{CO_2,DEB}(t)|\right) \label{eq:gamma}
\end{equation}

con $\kappa$ un hiperpar\'{a}metro de sensibilidad calibrado durante el entrenamiento.

\subsection{Dise\~{n}o Experimental}

\subsubsection{Material Biol\'{o}gico y Condiciones de Crianza}

BSF eggs were obtained from a laboratory colony maintained at Universidad Nacional de Colombia, Palmira. Larvae were reared in controlled climate chambers (Memmert HPP260, Germany) at $28 \pm 1^{\circ}$C and $65 \pm 3$\% relative humidity, with a 12h:12h light:dark photoperiod. Two substrate treatments were established:

Los huevos de BSF se obtuvieron de una colonia de laboratorio mantenida en la Universidad Nacional de Colombia, Palmira. Las larvas se criaron en c\'{a}maras clim\'{a}ticas controladas (Memmert HPP260, Alemania) a $28 \pm 1^{\circ}$C y $65 \pm 3$\% de humedad relativa, con fotoper\'{i}odo 12h:12h luz:oscuridad. Se establecieron dos tratamientos de sustrato:

\begin{itemize}
    \item \textbf{Alta calidad (AC):} Alimento comercial para pollos (20\% prote\'{i}na, 5\% l\'{i}pidos, 8\% cenizas);
    \item \textbf{Baja calidad (BC):} 50\% alimento para pollos + 50\% lodo anaer\'{o}bico desgasificado (12\% prote\'{i}na, 3\% l\'{i}pidos, 15\% cenizas).
\end{itemize}

La humedad del sustrato se mantuvo en $60 \pm 5$\% durante el per\'{i}odo de crianza de 20 d\'{i}as. Cada tratamiento comprendi\'{o} 3 r\'{e}plicas independientes (n=3), con densidad inicial de 1000 individuos por kg de sustrato h\'{u}medo.

\subsubsection{Protocolo de Monitoreo de CO$_2$}

Se construy\'{o} un sistema de respirometr\'{i}a de flujo cerrado utilizando:
\begin{itemize}
    \item Analizador infrarrojo de CO$_2$ WINSENT 410 (0--5000 ppm, resoluci\'{o}n 0,01\%);
    \item Controladores de flujo m\'{a}sico (Alicat MC Series, 0,1 L/min);
    \item C\'{a}maras herm\'{e}ticas de acr\'{i}lico de 2 L con sellos estancos.
\end{itemize}

El CO$_2$ se midi\'{o} a intervalos de 30 minutos (48 mediciones/d\'{i}a) durante 20 d\'{i}as, generando 960 puntos temporales por r\'{e}plica. Se rest\'{o} el CO$_2$ atmosf\'{e}rico basal (400 ppm) para obtener el CO$_2$ respiratorio neto.

\subsubsection{Cuantificaci\'{o}n de Biomasa}

El muestreo destructivo ocurri\'{o} en los d\'{i}as 5, 10, 15 y 20. En cada punto temporal, se muestrearon aleatoriamente 50 larvas, se ayunaron por 6 horas para vaciar el contenido intestinal, se enjuagaron con agua destilada y se secaron a 60$^{\circ}$C durante 48 horas hasta peso constante. El peso seco ($X_{DW}$) se registr\'{o} con precisi\'{o}n de 0,1 mg (Sartorius CPA225D).

\subsection{Calibraci\'{o}n y Validaci\'{o}n del Modelo}

\subsubsection{Estimaci\'{o}n de Par\'{a}metros DEB}

Los par\'{a}metros DEB ($a_{max}$, $B_{max,0}$, $\alpha$, $\beta$) se estimaron por minimizaci\'{o}n de cuadrados no lineales del Error Medio Absoluto (AME):

\begin{equation}
    AME = \frac{1}{n}\sum_{i=1}^{n} |X_{DW,obs,i} - X_{DW,modelo,i}| \label{eq:ame}
\end{equation}

utilizando el algoritmo Solver de Excel seg\'{u}n describe \citet{eriksenDynamicModellingFeed2022}. La tasa de mantenimiento $m = 0,08$ d\'{i}a$^{-1}$ y el costo de crecimiento $Y_B = 0,44$ se fijaron de valores de literatura \citep{laganaroGrowthMetabolicPerformance2021a}.

\subsubsection{Protocolo de Entrenamiento de la RNA-B}

La RNA-B se implement\'{o} en PyTorch 1.12. El entrenamiento emple\'{o}:
\begin{itemize}
    \item \textbf{Optimizador:} Adam con tasa de aprendizaje $10^{-3}$;
    \item \textbf{Funci\'{o}n de p\'{e}rdida:} P\'{e}rdida Huber (robusta a valores at\'{i}picos);
    \item \textbf{Regularizaci\'{o}n:} Decaimiento de peso L2 ($\lambda = 10^{-4}$);
    \item \textbf{Validaci\'{o}n:} Validaci\'{o}n cruzada dejando-una-r\'{e}plica-fuera (5 pliegues);
    \item \textbf{Parada temprana:} Paciencia de 50 \'{e}pocas con monitoreo de p\'{e}rdida de validaci\'{o}n.
\end{itemize}

Las caracter\'{i}sticas de entrada se estandarizaron (Z-score: $\mu = 0$, $\sigma = 1$) para garantizar estabilidad num\'{e}rica.

\subsubsection{Modelos de Referencia}

El modelo h\'{i}brido se compar\'{o} contra:
\begin{enumerate}
    \item \textbf{DEB puro:} Modelo de Eriksen (2022) con par\'{a}metros de literatura;
    \item \textbf{RNA-B pura:} Red neuronal con arquitectura id\'{e}ntica pero sin precondicionamiento DEB;
    \item \textbf{Regresi\'{o}n log\'{i}stica:} Ajuste cl\'{a}sico de curva de crecimiento.
\end{enumerate}

Las m\'{e}tricas de desempe\~{n}o incluyeron AME, Ra\'{i}z del Error Cuadr\'{a}tico Medio (RMSE) y coeficiente de determinaci\'{o}n ($R^2$).

\section{Resultados}

\subsection{Trayectorias de CO$_2$ y Efectos del Sustrato}

[Espacio reservado para figuras: curvas temporales de CO$_2$ para tratamientos AC vs BC]

\subsection{Desempe\~{n}o del Modelo DEB}

[Espacio reservado: resultados de ajuste DEB, tablas de par\'{a}metros, valores AME]

\subsection{Eficacia de la Correcci\'{o}n RNA-B}

[Espacio reservado: an\'{a}lisis de residuos, importancia de caracter\'{i}sticas, pesos de la red]

\subsection{Validaci\'{o}n del Modelo H\'{i}brido}

[Espacio reservado: resultados de validaci\'{o}n cruzada, tabla comparativa, intervalos de predicci\'{o}n]

\section{Discusi\'{o}n}

\subsection{Interpretabilidad Fisiol\'{o}gica}

El marco h\'{i}brido preserva la transparencia mecanicista del modelado DEB mientras acomoda las desviaciones metab\'{o}licas espec\'{i}ficas del sustrato. El t\'{e}rmino de correcci\'{o}n $\gamma(t) \cdot \hat{B}_{RNA}(t)$ puede interpretarse como un \textit{factor de estr\'{e}s ambiental} no capturado por la formulaci\'{o}n DEB est\'{a}ndar -- an\'{a}logo a los coeficientes de modulaci\'{o}n $\alpha$ y $\beta$ en \citet{eriksenDynamicModellingFeed2022}, pero aprendido adaptativamente de los datos en lugar de ajustado manualmente.

\subsection{Escalabilidad Industrial}

Para instalaciones industriales de BSF, el modelo h\'{i}brido permite:
\begin{itemize}
    \item \textbf{Monitoreo de biomasa en tiempo real:} Los sensores de CO$_2$ cuestan US\$50--200 vs. US\$500+ para sistemas de visi\'{o}n;
    \item \textbf{Optimizaci\'{o}n de cosecha:} Predecir la llegada de $B_{max}$ para programar la recolecci\'{o}n de prepupas;
    \item \textbf{Contabilidad de GEI:} Cuantificaci\'{o}n directa de emisiones de CO$_2$ por kg de biomasa producida.
\end{itemize}

\subsection{Limitaciones y Trabajo Futuro}

Las limitaciones actuales incluyen:
\begin{enumerate}
    \item La replicaci\'{o}n limitada (n=5) restringe la generalizaci\'{o}n a sustratos novedosos;
    \item La suposici\'{o}n de constancia de temperatura-humedad; el trabajo futuro debe integrar covariables ambientales;
    \item Se asumi\'{o} CH$_4$ despreciable; los sustratos metanog\'{e}nicos pueden requerir modelado de gas dual.
\end{enumerate}

\section{Conclusiones}

Este estudio demuestra que un marco DEB-RNA h\'{i}brido puede estimar con precisi\'{o}n la biomasa larvaria de BSF a partir de mediciones respiratorias de CO$_2$, superando tanto los enfoques puramente mecanicistas como puramente basados en datos. El modelo proporciona una soluci\'{o}n escalable y fundamentada fisiol\'{o}gicamente para el monitoreo no destructivo de biomasa en agricultura de precisi\'{o}n con insectos, con aplicaciones directas para la optimizaci\'{o}n de procesos y la evaluaci\'{o}n de impacto ambiental en sistemas de bioeconom\'{i}a circular.

\section*{Agradecimientos}

Esta investigaci\'{o}n fue financiada por [n\'{u}mero de convocatoria MinCiencias/Colombia]. Los autores agradecen a [t\'{e}cnicos de laboratorio] por el manejo animal y a [proveedor de sensores] por el pr\'{e}stamo de equipos.

\section*{Referencias}

\begin{thebibliography}{20}

\bibitem[Arrese y Soulages(2010)]{arreseInsectFatBody2010}
Arrese, E.L., Soulages, J.L., 2010. Insect fat body: energy, metabolism, and regulation. \textit{Annu. Rev. Entomol.} 55, 207--225.

\bibitem[Diener et al.(2009)]{dienerConversionOrganicMaterial2011}
Diener, S., Zurbr\'{u}gg, C., Tockner, K., 2009. Conversion of organic material by black soldier fly larvae: establishing optimal feeding rates. \textit{Waste Manage. Res.} 27, 603--610.

\bibitem[Eriksen(2022)]{eriksenDynamicModellingFeed2022}
Eriksen, N.T., 2022. Dynamic modelling of feed assimilation, growth, lipid accumulation, and CO$_2$ production in black soldier fly larvae. \textit{PLoS ONE} 17(10), e0276605.

\bibitem[Gold et al.(2020)]{goldBiowasteTreatmentBlack2020}
Gold, M., Cassar, C.M., Zurbr\'{u}gg, C., Kreuzer, M., Boulos, S., Diener, S., et al., 2020. Biowaste treatment with black soldier fly larvae: increasing performance through the formulation of biowastes based on protein and carbohydrates. \textit{Waste Manage.} 102, 319--329.

\bibitem[Kok(2021)]{kokPreliminaryProjectDesign2021}
Kok, R., 2021. Preliminary project design for insect production: part 1 -- overall mass and energy/heat balances. \textit{J. Insects Food Feed} 7, 499--509.

\bibitem[Laganaro et al.(2021)]{laganaroGrowthMetabolicPerformance2021a}
Laganaro, M., Bahrndorff, S., Eriksen, N.T., 2021. Growth and metabolic performance of black soldier fly larvae grown on low and high-quality substrates. \textit{Waste Manage.} 121, 198--205.

\bibitem[LeCun et al.(2015)]{lecunDeepLearning2015}
LeCun, Y., Bengio, Y., Hinton, G., 2015. Deep learning. \textit{Nature} 521, 436--444.

\bibitem[Li et al.(2015)]{liSimultaneousUtilizationGlucose2015}
Li, W., Li, M., Zheng, L., Liu, Y., Zhang, Y., Yu, Z., et al., 2015. Simultaneous utilization of glucose and xylose for lipid accumulation in black soldier fly. \textit{Biotechnol. Biofuels} 8, 11.

\bibitem[Makkar et al.(2014)]{makkarStateoftheartUseInsects2014}
Makkar, H.P.S., Tran, G., Heuz\'{e}, V., Ankers, P., 2014. State-of-the-art on use of insects as animal feed. \textit{Anim. Feed Sci. Technol.} 197, 1--33.

\bibitem[Meneguz et al.(2018)]{meneguzEffectRearingSubstrate2018}
Meneguz, M., Schiavone, A., Gai, F., Dama, A., Lussiana, C., Renna, M., et al., 2018. Effect of rearing substrate on growth performance, waste reduction efficiency and chemical composition of black soldier fly larvae. \textit{J. Sci. Food Agric.} 98, 5776--5784.

\bibitem[Padmanabha et al.(2020)]{padmanabhaComprehensiveDynamicGrowth2020}
Padmanabha, M., Kobelski, A., Hempel, A.-J., Streif, S., 2020. A comprehensive dynamic growth and development model of \textit{Hermetia illucens} larvae. \textit{PLoS ONE} 15, e0239084.

\bibitem[Prasetya et al.(2021)]{prasetya2021growth}
Prasetya, A., Darmawan, R., Araujo, T.L.B., Petrus, H.T.B.M., Setiawan, F.A., 2021. A growth kinetics model for black soldier fly larvae. \textit{Int. J. Technol.} 12, 207--216.

\bibitem[Sripontan et al.(2020)]{sripontan2020modeling}
Sripontan, Y., Chiu, C.-I., Tanansathaporn, S., Leasen, K., Manlong, K., 2020. Modeling the growth of black soldier fly \textit{Hermetia illucens}: an approach to evaluate diet quality. \textit{J. Econ. Entomol.} 113, 742--751.

\bibitem[Johnson et al.(2019)]{kamilarisDeepLearningAgriculture2018}
Johnson, K.B., et al., 2019. Deep learning in agriculture: a survey. \textit{Comput. Electron. Agric.} 163, 104859.

\end{thebibliography}

\end{document}
